\input{../tex-inputs/latex-leseansicht-vorspann}

\section[Arthur Schnitzler an Elsa Plessner, 17. 1. 1916]{L04340 Arthur Schnitzler an Elsa Plessner, 17. 1. 1916}
\nopagebreak\mylabel{L04340v}
\rehead{ }\normalsize\beginnumbering\briefempfaengerindex{Plessner, Elsa@\textsc{Plessner, Elsa}!zzzSchnitzler, Arthur@\emph{von Arthur Schnitzler}!1916-01-171@{17. 1. 1916}|(be}
\toendnotes[C]{\smallbreak\pagebreak[2]}
\correspDesc{Versand  durch Arthur Schnitzler am 17. 1. 1916 in Wien
\newline{}Erhalt  durch Elsa Plessner im Zeitraum [18. 1. 1916
                  – 22. 1. 1916?] in München}\toendnotes[C]{\smallbreak}
\Standort{DLA, A:Schnitzler, HS.1985.1.829.}
\physDesc{Brief, 3 Blätter, 5 Seiten, 3890 Zeichen
\newline{}Schreibmaschine
\newline{}Handschrift Schreibkraft: Bleistift (\noindent{}Ergänzungen von Satzzeichen)
\newline{}Handschrift Arthur Schnitzler: 1) roter Buntstift, lateinische Kurrent (\noindent{}Vermerk »Ginsberg« und »\textcolor{gray}{Pl}« auf jedem Blatt und eine Unterstreichung)\hspace{1em}2) Bleistift, lateinische Kurrent (\noindent{}Ergänzung der Datumsangabe auf Blatt 2 und 3: »17/1 916«)\hspace{1em}}\toendnotes[C]{\smallbreak}
\pstart
           \raggedleft{}{\pb}17. 1. 1916.\pend
           
\pstart\center{}Verehrte gnädige Frau.\pend\vspace{0.5em}
\pstart
           Ihrem freundlichen Wunsch entsprechend habe ich das erste Couvert eröffnet und Ihren
                  \label{K_L04340-1v}\edtext{Brief}{\lemma{\textnormal{\emph{Brief}}}\Cendnote{\textnormal{XXXX Auszeichnungsfehler: Dokument L03730 nicht gefunden.}}}\label{K_L04340-1}
               gelesen. Die darin enthaltene Entstehungsanalyse Ihres Stücks\pwindex{Plessner, Elsa 22.\,8.\,1875 Wien – 7.\,5.\,1932 Alicante@\textsc{Plessner, Elsa} (22.\,8.\,1875 Wien – 7.\,5.\,1932 Alicante), \emph{Schriftstellerin, Schauspielerin}!Musik@\strich\emph{Musik}|pwv} hat mich so weit angeregt, dass ich
               von meinem Prinzip abgehend mich an die Lektüre\pwindex{Plessner, Elsa 22.\,8.\,1875 Wien – 7.\,5.\,1932 Alicante@\textsc{Plessner, Elsa} (22.\,8.\,1875 Wien – 7.\,5.\,1932 Alicante), \emph{Schriftstellerin, Schauspielerin}!Musik@\strich\emph{Musik}|pwv} herangewagt habe. Nachdem ich sie beendet kann ich
               Ihnen meine Empfindung nicht verschweigen, dass Sie in der Ausführung um ein
               Beträchtliches hinter Ihren Absichten zurückgeblieben sind. Was Ihnen als Stil
               erscheint, stellt sich meines Erachtens eher als eine gewisse Kühle, ja in den
               entscheidenden Momenten Dürftigkeit der Diktion dar und wo es auf das Wort ankäme
               sehe ich nicht viel mehr (zuweilen) als Pantomine. Sie könnten mir vielleicht
               antworten, dass gerade in den grossen entscheidenden Momenten des Daseins das Wort
               notwendig versagt; aber wäre ich selbst {\pb}geneigt eine solche Entschuldigung auch im Drama gelten zu lassen, so müsste ich
               zumindest verlangen, dass mir die Gestalt, der das dramatische Wort nicht gegeben
               ist, mir vorher interessant und lebendig genug wurde, um mir auch noch in ihrem
               Schweigen oder in der Schwächlichkeit ihrer Ausdrucksweise interessant, lebendig –
               und verständlich zu sein. Aber Ihre Gestalten auch Reiter\pwindex{Plessner, Elsa 22.\,8.\,1875 Wien – 7.\,5.\,1932 Alicante@\textsc{Plessner, Elsa} (22.\,8.\,1875 Wien – 7.\,5.\,1932 Alicante), \emph{Schriftstellerin, Schauspielerin}!Musik@\strich\emph{Musik}|pwv} und Jolanthe\pwindex{Plessner, Elsa 22.\,8.\,1875 Wien – 7.\,5.\,1932 Alicante@\textsc{Plessner, Elsa} (22.\,8.\,1875 Wien – 7.\,5.\,1932 Alicante), \emph{Schriftstellerin, Schauspielerin}!Musik@\strich\emph{Musik}|pwv} sind recht blass geblieben. Es fehlt das
               charakteristische Detail, das geheimnisvoll Atmosphärische, jenes Undefinierbare, das
               der erdichteten Figur innerhalb des Kunstwerks Vergangenheit, Zukunft und damit erst
               die richtige Gegenwart verleiht. Warum glaube ich nicht an die Liebe auf den ersten
               Blick zwischen Jolanthe\pwindex{Plessner, Elsa 22.\,8.\,1875 Wien – 7.\,5.\,1932 Alicante@\textsc{Plessner, Elsa} (22.\,8.\,1875 Wien – 7.\,5.\,1932 Alicante), \emph{Schriftstellerin, Schauspielerin}!Musik@\strich\emph{Musik}|pwv} und
                  Reiter\pwindex{Plessner, Elsa 22.\,8.\,1875 Wien – 7.\,5.\,1932 Alicante@\textsc{Plessner, Elsa} (22.\,8.\,1875 Wien – 7.\,5.\,1932 Alicante), \emph{Schriftstellerin, Schauspielerin}!Musik@\strich\emph{Musik}|pwv}? Weil ich an sie
               selber nicht glaube. Aber nehme ich selbst diese Tatsache, was mir ohne Liebestrank
               und ohne musikalische Begleitung schwer genug fiele, als gegeben hin, so wehre ich
               mich durchaus gegen alles, was folgt als theaterhaft {\pb}und jeder Wahrheit bar. Hier ist nicht nur
               äussere Unwahrscheinlichkeit, wogegen sich auch schon genug einwenden liesse, hier
               ist ein Mangel an innerer Wahrheit, den durch die Gesetze theatralischer Verkürzung,
               auch nicht durch den besonderen Stil, den Sie hier vielleicht gefunden zu haben
               glauben, gerechtfertigt erscheint. Das ganze Gebahren der Jolanthe\pwindex{Plessner, Elsa 22.\,8.\,1875 Wien – 7.\,5.\,1932 Alicante@\textsc{Plessner, Elsa} (22.\,8.\,1875 Wien – 7.\,5.\,1932 Alicante), \emph{Schriftstellerin, Schauspielerin}!Musik@\strich\emph{Musik}|pwv} wäre durch eine pathologische
               Disposition zu erklären; diese ist aber vorher weder angedeutet, noch vorbereitet und
               so sehen wir uns der Willkür des Autors hilflos ausgeliefert, was als eine
               Hauptquelle des ästhetischen Unbehagens gelten darf. All dem gegenüber fallen die
               kleineren Peinlichkeiten kaum mehr ins Gewicht, die gewissermassen nur aus der
               ungenügenden Beherrschung des Stoffes entspringen und milder als scenische
               Ungeschichlichkeiten anzusprechen wären. Dass Jolanthe\pwindex{Plessner, Elsa 22.\,8.\,1875 Wien – 7.\,5.\,1932 Alicante@\textsc{Plessner, Elsa} (22.\,8.\,1875 Wien – 7.\,5.\,1932 Alicante), \emph{Schriftstellerin, Schauspielerin}!Musik@\strich\emph{Musik}|pwv} am Ende durch ihren Selbstmordversuch erblindet,
               wirkt keineswegs tragisch, sondern als ein unglück{\pb}seliger Zufall, wie sich ihn wohl
                  de\textcolor{gray}{r}{ }\damage{liebe{ }\textcolor{gray}{Gott}} erlauben darf, aber niemals der dramatische Autor. Die Musik in Ihrem Stück\pwindex{Plessner, Elsa 22.\,8.\,1875 Wien – 7.\,5.\,1932 Alicante@\textsc{Plessner, Elsa} (22.\,8.\,1875 Wien – 7.\,5.\,1932 Alicante), \emph{Schriftstellerin, Schauspielerin}!Musik@\strich\emph{Musik}|pwv} vermag ich nicht zu
               spüren. Ob der Fehler an Ihnen liegt oder an mir müssen wir dahingestellt sein
               lassen; mir bleibt am Ende nichts übrig, auf die Gefahr des Irrtums hin, als Ihnen
               von meinem rein persönlichen Eindruck zu berichten, der sich vor
               allem auf das Ganze
               Ihrer Arbeit bezieht. Wenn ich Ihnen zum Schluss noch sage, dass ich im Dialog da und
               dort manche feine und kluge Bemerkung gefunden habe, und dass der technische Aufbau
               Ihres dramatischen Gerüstes eine keineswegs ungeübte Hand, dass endlich die
               Problemstellung als solche eine gewisse theatralische und sogar seelische Courage
               verrät, so ist Ihnen nicht viel damit geholfen, aber ich möchte doch auch nicht, dass
               Sie als Höflichkeitsphrasen nehmen, was geradeso ehrlich gemeint ist, wie das minder
               Höfliche, was ich Ihnen im ersten Teil meines Briefes vorgebracht habe.\pend
           
\pstart
           {\pb}Ihrem zweiten \label{K_L04340-2v}\edtext{Briefe}{\lemma{\textnormal{\emph{Briefe}}}\Cendnote{\textnormal{XXXX Auszeichnungsfehler: Dokument L03731 nicht gefunden.}}}\label{K_L04340-2} lag keine Postkarte bei.
               Ich erwarte nun eine Weisung von Ihnen wohin ich Ihr Manuscript\pwindex{Plessner, Elsa 22.\,8.\,1875 Wien – 7.\,5.\,1932 Alicante@\textsc{Plessner, Elsa} (22.\,8.\,1875 Wien – 7.\,5.\,1932 Alicante), \emph{Schriftstellerin, Schauspielerin}!Musik@\strich\emph{Musik}|pwv} zu senden habe.\pend
           
\pstart
           Mit verbindlichsten Grüssen{\\[\baselineskip]}Ihr Sie aufrichtig schätzender und keineswegs
               unfehlbarer\pend
           \leftskip=0em{}\selectlanguage{ngerman}\endnumbering\briefempfaengerindex{Plessner, Elsa@\textsc{Plessner, Elsa}!zzzSchnitzler, Arthur@\emph{von Arthur Schnitzler}!1916-01-171@{17. 1. 1916}|)be}\mylabel{L04340h}  \newcommand{\dateiname}{L04340}\newcommand{\titel}{Arthur Schnitzler an Elsa Plessner, 17. 1. 1916}\newcommand{\editorInnen}{Herausgegeben von Jahnke, SelmaMüller, Martin Anton}\input{../tex-inputs/latex-leseansicht-abspann}
